\documentclass[UTF8,11pt,a4paper]{ctexart}
\usepackage{amsmath,amssymb}
\usepackage[margin=1in]{geometry}
\usepackage{booktabs}
\usepackage{enumitem}
\usepackage[hidelinks]{hyperref}
\usepackage[expansion=false,protrusion=true]{microtype}

\newcommand{\R}{\mathbb{R}}

\title{\textbf{ECT 项目交接报告 — 2026-08-10(Role C)}%
\thanks{面向项目 leader 与 mentor 的汇报。对应仓库 `recurrence\_of\_ect`,分支
`role-c/g13-vs-g10-fid-0809`。}}
\author{Role C --- 理论负责人}
\date{2026-08-10}

\begin{document}
\maketitle

\begin{abstract}
本报告总结 2026-08-10 围绕 ICLR 2027 战略的实测工作。核心结论:通过一系列
\emph{实测}实验,把方案的生死判断钉死——非标量梯度残差存在、有结构、状态条件、
普遍、\emph{良性},但\emph{附庸于标量 gap}。因此"机制+方法"与"诊断信号"两条
升级路径均被证伪,论文最终定位为 \emph{workshop 级诚实负结果},不投 ICLR 主会。
\end{abstract}

\tableofcontents

% ------------------------------------------------------------------
\section{战略轨迹}
\label{sec:trajectory}

\begin{center}
\begin{tabular}{@{}ll@{}}
\toprule
阶段 & 结论 \\
\midrule
机制+方法($40$--$50\%$) & 被 $g{=}1.3$ 赢 FID 证伪 \\
诚实诊断($20$--$30\%$) & 被"良性"非结果质疑 \\
诊断信号转向($30$--$35\%$) & 被 P1 附庸性检验证伪(偏相关 $\approx 0$) \\
\bottomrule
\end{tabular}
\end{center}

最终定位:\textbf{workshop 级诚实负结果}。

% ------------------------------------------------------------------
\section{实测诊断证据}
\label{sec:evidence}

六个实验,全部实测:

\begin{center}
\begin{tabular}{@{}lll@{}}
\toprule
环节 & 实验 & 结论 \\
\midrule
存在 & moment-memory & 标量 null 被证伪($h$ $1.001$ vs.\ $0.837$, $\mathrm{Corr}\approx 0$) \\
有结构 & E5 & 深度趋势($\rho\approx 0.4$)+ 剂量响应,非噪声 \\
无害 & E6 + 干净 FID & 与 FID 反相关 $r{=}{-}0.99$;$g{=}1.3$ 赢 $-34\%$/$-45\%$ \\
状态条件 & E7 & $R_{\mathrm{grad}}$ 随 $n_K$ 增长($0.031\to0.091$) \\
稳健 & E11 & 跨 3 seed(1024k 5/6)+ 支撑阈值 \\
普遍 & E3 & 梯度残差跨 optimizer 复现(AdamW $0.095$ vs.\ RAdam $0.091$) \\
\bottomrule
\end{tabular}
\end{center}

\textbf{关键数字:} $a_K^*{=}0.761$、$R_{\mathrm{grad}}{=}5.85\%$、
$R_{\mathrm{opt}}{=}8.57\%$、$H_K{=}R_{\mathrm{opt}}$ 精确($1.7\times10^{-18}$)、
$h_{\mathrm{pred}}{=}1.001$ vs.\ $h_{\mathrm{actual}}{=}0.837$。

% ------------------------------------------------------------------
\section{P1 决定性结果}
\label{sec:p1}

\textbf{残差附庸于标量 gap。} 在 7 个 gap 上实测残差,6 个共同 gap 上计算:

\begin{center}
\begin{tabular}{@{}ll@{}}
\toprule
相关 & 值 \\
\midrule
$\mathrm{corr}(\text{残差},\text{FID})$ & $-0.817$ \\
$\mathrm{corr}(|g-1|,\text{FID})$ & $-0.899$ \\
$\mathrm{corr}(\text{残差},|g-1|)$ & $+0.938$ \\
\textbf{偏相关(残差, FID $|$ $|g-1|$)} & \textbf{$+0.168\approx 0$} \\
\bottomrule
\end{tabular}
\end{center}

残差几乎由 $|g-1|$ 决定,在标量 gap 之外\textbf{不增加对 FID 的预测力}。
诊断信号论题不成立。

% ------------------------------------------------------------------
\section{最终定位}
\label{sec:position}

\begin{quote}
非标量梯度残差存在、有结构、状态条件、普遍、\emph{良性}——但\emph{附庸于标量
gap}。它是少步训练的真实特征,但不是质量诊断信号,也不有害。\textbf{论文定位:
workshop 级诚实负结果,不投 ICLR 主会。}
\end{quote}

% ------------------------------------------------------------------
\section{交付物}
\label{sec:deliverables}

\begin{itemize}
  \item \textbf{诊断论文}: `docs/ICLR2027\_DIAGNOSIS\_PAPER.{md,tex,pdf}`
  \item \textbf{战略文档}: `docs/ICLR2027\_STRATEGY.{tex,pdf,md}`
  \item \textbf{对抗性审查}: `docs/ICLR2027\_PLAN\_REVIEW.md`
  \item \textbf{实验结果}(7 份): `analysis/g13\_vs\_g10\_fid\_result.md`、
        `e5/e6/e7/e11/e3/p1\_*\_result.md`
  \item \textbf{Git}: 分支 `role-c/g13-vs-g10-fid-0809`,23 个 commit 领先 main,
        未合并 main
\end{itemize}

% ------------------------------------------------------------------
\section{待办}
\label{sec:todo}

\begin{enumerate}[leftmargin=*,itemsep=1pt]
  \item 本地 push 分支(勿推 main)
  \item 把论文定位更新为 workshop 级诚实负结果
  \item 可选:投 ``I Can't Believe It's Not Better'' workshop
\end{enumerate}

\end{document}
